% !TeX root = ../main.tex

\sysusetup{
  keywords  = {多无人机，移动边缘计算，轨迹优化，任务卸载，Attention-MAPPO，多智能体强化学习，双层优化，Lagrange 约束},
  keywords* = {Multi-UAV, mobile edge computing, trajectory optimization, task offloading, Attention-MAPPO, multi-agent reinforcement learning, hierarchical optimization, Lagrange constraints},
}

\begin{abstract}
多无人机辅助移动边缘计算能够通过空中节点的机动部署提升边缘覆盖、任务接入和协同计算能力，但系统性能同时受到无人机轨迹、无线链路状态、请求级卸载决策、截止期约束和宏基站回传负载的共同影响。若将无人机连续轨迹控制与多请求离散卸载决策直接合并为单一联合动作空间，训练复杂度和在线决策开销都会迅速上升。为此，本文提出一种基于 Attention-MAPPO 的多无人机 MEC 轨迹与任务卸载协同优化框架。该框架将协同调度拆分为上层轨迹控制和下层请求级卸载两个决策层：上层采用 attention-MAPPO 建模无人机之间的协同关系并输出轨迹动作；下层采用 constrained attention offload MAPPO，在每个有效服务请求上选择本地无人机执行、协作无人机执行或宏基站执行。本文进一步设计质量感知动作 mask 屏蔽明显不可行的协作与宏基站动作，并在下层奖励中引入 Lagrange 约束项刻画截止期满足率与宏基站负载之间的权衡。同时，本文在系统模型中显式建模无人机飞行能耗和用户设备电池动态，并通过无线能量传输机制维持终端设备在线。实验基于 5 架 UAV、100 个用户设备、3 个 training seeds 和 10 组 workload seeds 展开。结果表明，完整双层 MARL 相比 uncoordinated greedy + heuristic 基线可将平均 reward 从 -5933.9 提升到 -1451.7（$\Delta = +4482.2$, $p = 1.0\times10^{-32}$），UAV 侧能耗从 114.70M 降至 58.06M（$\Delta = -56.64M$, $p = 1.7\times10^{-19}$），公平性由 0.7745 提升至 0.9363（$\Delta = +0.162$, $p = 1.7\times10^{-11}$），offline rate 由 0.58\% 降至 0.00\%。消融实验表明，Lagrange 约束、质量 mask 和 attention 模块会显著改变卸载分布与能耗/负载权衡。进一步地，DSR-priority 配置通过调整 reward 权重与 Lagrange 约束阈值，验证了该框架在不同服务质量偏好下的可调节性。
\end{abstract}

\begin{abstract*}
Multi-UAV assisted mobile edge computing (MEC) enhances edge coverage, task admission, and cooperative computing through aerial node deployment, yet system performance is jointly affected by UAV trajectory, wireless link dynamics, per-request offloading decisions, deadline constraints, and macro base station (MBS) backhaul load. Directly combining continuous UAV trajectory control with discrete per-request offloading decisions into a single joint action space leads to prohibitive training complexity and online decision latency. This paper proposes an Attention-MAPPO based collaborative trajectory and task offloading optimization framework for multi-UAV MEC. The framework decomposes coordinated scheduling into an upper-layer trajectory control module and a lower-layer per-request offloading module. The upper layer employs attention-MAPPO to model inter-UAV coordination and output trajectory actions; the lower layer employs constrained attention offload MAPPO to select, for each service request, among local UAV execution, cooperative UAV execution, and MBS execution. We design a quality-aware action mask that prunes clearly infeasible cooperative and MBS actions, and introduce Lagrange constraint terms in the lower-layer reward to characterize the trade-off between deadline satisfaction rate (DSR) and MBS load ratio. Additionally, we explicitly model UAV flight energy consumption and UE battery dynamics, maintaining terminal device availability via wireless power transfer (WPT). Experiments are conducted across 3 training seeds and 10 workload seeds on a system with 5 UAVs and 100 UEs. Results show that the full hierarchical MARL improves mean reward from -5933.9 to -1451.7 relative to the uncoordinated greedy + heuristic baseline ($\Delta = +4482.2$, $p = 1.0\times10^{-32}$), reduces UAV-side energy from 114.70M to 58.06M ($\Delta = -56.64M$, $p = 1.7\times10^{-19}$), raises fairness from 0.7745 to 0.9363 ($\Delta = +0.162$, $p = 1.7\times10^{-11}$), and lowers the offline rate from 0.58\% to 0.00\%. Ablation studies demonstrate that Lagrange constraints, quality masks, and attention modules significantly alter offloading distribution and energy-load trade-offs. Furthermore, the DSR-priority configuration validates the framework's adjustability across different service quality preferences.
\end{abstract*}
